\documentclass[a4paper,12pt]{article}

\usepackage[utf8]{inputenc}
\usepackage[portuguese]{babel}
\usepackage{fancyhdr}
\usepackage{amsmath}
\usepackage{amsthm}
\usepackage{amssymb}
\usepackage[portuguese,algoruled,lined,linesnumbered]{algorithm2e}
\usepackage{tikz}
\usetikzlibrary{positioning,shadows,arrows.meta,calc}

% ---- PREENCHA ----
\newcommand{\Lista}{6}
\newcommand{\Nome}{Leonardo Heidi Almeida Murakami}
\newcommand{\NUSP}{11260186}
% -----------------------------

\pagestyle{fancy}
\fancyhf{}
\fancyhead[L]{\Nome}
\fancyhead[R]{NUSP \NUSP}
\fancyfoot[L]{Lista \Lista}
\fancyfoot[C]{MAC0338}

\begin{document}
\section*{Exercício 1}
\subsection*{Menor s[i]}
Exemplo que prova que o algoritmo guloso que escolhe o menor s[i] não é ótimo:

\begin{center}
\begin{tikzpicture}[x=2.0cm, y=0.8cm, thick]
  % timeline
  \draw[->] (0,0) -- (5.5,0) node[below right] {$t$};
  \foreach \x in {1,2,3,4,5} {
    \draw (\x,0.15) -- (\x,-0.15) node[below] {\x};
  }

  % tarefas
  % i=1: t=1 até t=5
  \draw[fill=blue!20] (1,0.3) rectangle (5,0.8);
  \node[left] at (1,0.55) {$i=1$};

  % i=2: t=2 até t=3
  \draw[fill=red!20] (2,1.0) rectangle (3,1.5);
  \node[left] at (2,1.25) {$i=2$};

  % i=3: t=3 até t=4
  \draw[fill=green!20] (3,1.7) rectangle (4,2.2);
  \node[left] at (3,1.95) {$i=3$};

  % i=4: t=4 até t=5
  \draw[fill=yellow!40] (4,2.4) rectangle (5,2.9);
  \node[left] at (4,2.65) {$i=4$};

\end{tikzpicture}
\end{center}
Neste caso escolheriamos a tarefa 1, com $S_1 = \left[1\right]$ enquanto a solucao otima é $S_o = \left[2, 3, 4\right]$.

\subsection*{Menor f[i] - s[i]}
Exemplo que prova que o algoritmo guloso que escolhe o menor $f[i] - s[i]$ (tempo de duração mais curto da tarefa) não é ótimo:

\begin{center}
\begin{tikzpicture}[x=1.2cm, y=0.8cm, thick]
  % timeline
  \draw[->] (0,0) -- (9,0) node[below right] {$t$};
  \foreach \x in {1,2,3,4,5,6,7,8} {
    \draw (\x,0.15) -- (\x,-0.15) node[below] {\x};
  }

  % tarefas
  % i=1: t=1 até t=4
  \draw[fill=blue!20] (1,0.3) rectangle (4,0.8);
  \node[left] at (1,0.55) {$i=1$};

  % i=2: t=4 até t=8
  \draw[fill=red!20] (4,1.0) rectangle (8,1.5);
  \node[left] at (4,1.25) {$i=2$};

  % i=3: t=3 até t=5
  \draw[fill=green!20] (3,1.7) rectangle (5,2.2);
  \node[left] at (3,1.95) {$i=3$};

\end{tikzpicture}
\end{center}

Se escolhermos sempre a tarefa de menor duração ($f[i] - s[i]$), escolheríamos primeiro $i=3$ (duração 3) e depois não poderíamos escolher as outras duas devido a sobreposição de intervalos. Mas a solução ótima seria selecionar as tarefas $i=2$ e $i=3$ (que não se sobrepõem), ou seja, $S_o = [1, 2]$.

\subsection*{Menor quantidade de conflitos}
\begin{center}
\begin{tikzpicture}[x=0.9cm, y=0.65cm, thick]
  % timeline
  \draw[->] (0,0) -- (11,0) node[below right] {$t$};
  \foreach \x in {1,2,...,10} {
    \draw (\x,0.15) -- (\x,-0.15) node[below] {\x};
  }

  % tarefas
  % i=1: t=1 até t=3
  \draw[fill=blue!20] (1,0.3) rectangle (3,0.8);
  \node[left] at (1,0.55) {$i=1$};

  % i=2: t=3 até t=5
  \draw[fill=red!20] (3,1.0) rectangle (5,1.5);
  \node[left] at (3,1.25) {$i=2$};

  % i=3: t=6 até t=8
  \draw[fill=green!20] (6,1.7) rectangle (8,2.2);
  \node[left] at (6,1.95) {$i=3$};

  % i=4: t=8 até t=10
  \draw[fill=yellow!40] (8,2.4) rectangle (10,2.9);
  \node[left] at (8,2.65) {$i=4$};

  % i=5: t=4 até t=7
  \draw[fill=orange!35] (4,3.1) rectangle (7,3.6);
  \node[left] at (4,3.35) {$i=5$};

  % i=6: t=2 até t=4
  \draw[fill=purple!20] (2,3.8) rectangle (4,4.3);
  \node[left] at (2,4.05) {$i=6$};

  % i=7: t=7 até t=9
  \draw[fill=pink!35] (7,4.5) rectangle (9,5.0);
  \node[left] at (7,4.75) {$i=7$};

  % i=8: t=2 até t=4
  \draw[fill=orange!30] (2,5.2) rectangle (4,5.7);
  \node[left] at (2,5.45) {$i=8$};

  % i=9: t=7 até t=9
  \draw[fill=teal!40] (7,5.9) rectangle (9,6.4);
  \node[left] at (7,6.15) {$i=9$};

  % i=10: t=2 até t=4
  \draw[fill=gray!25] (2,6.6) rectangle (4,7.1);
  \node[left] at (2,6.85) {$i=10$};

  % i=11: t=7 até t=9
  \draw[fill=brown!35] (7,7.3) rectangle (9,7.8);
  \node[left] at (7,7.55) {$i=11$};
\end{tikzpicture}
\end{center}
Primeiro escolhemos i=5 pela menor quantidade de conflitos e menor s[i], i=5 nao faz parte da solucao otima $S_o = \left[1, 2, 3, 4\right]$.

\section*{Exercício 2}
\section*{Exercício 3}
\section*{Exercício 4}
\section*{Exercício 5}
\section*{Exercício 6}
\section*{Exercício 7}
\section*{Exercício 8}
\end{document}